\section{Relational event models and case-control estimation}
\label{sec:background}

This section fixes the statistical objects that the rest of the paper
builds on. We recall the relational event intensity and its partial
likelihood (Section~\ref{sec:rem-likelihood}), the nested case-control
device that makes that likelihood computable at scale and its expression
as a conditional-logistic (risk-set softmax) partial likelihood
(Section~\ref{sec:case-control}), the algebraically related degenerate
case-1-control binomial reduction used by the smooth backend
(Section~\ref{sec:degenerate}), the endogenous-statistics catalogue that
supplies the predictors (Section~\ref{sec:statistics}), and the
time-shifted partial likelihood that admits global covariates
(Section~\ref{sec:global}). None of this machinery is new to \pkg{amore};
the contribution of the package is to consolidate it. We give just enough
notation to make the three estimation backends of
Section~\ref{sec:framework} and the goodness-of-fit family of
Section~\ref{sec:software-gof} precise.

\subsection{The relational event intensity and its partial likelihood}
\label{sec:rem-likelihood}

A relational event log is an ordered sequence of triples
$e_i = (s_i, r_i, t_i)$, $i = 1, \dots, M$, with $t_1 < t_2 < \cdots <
t_M$, where at time $t_i$ sender $s_i$ directs an action at receiver
$r_i$ \citep{butts2008}. Let $\mathcal{R}(t)$ denote the \emph{risk set}
at time $t$: the collection of ordered dyads $(s, r)$ that could fire an
event just after $t$. In the simplest one-mode setting with actor set
$V$, $\mathcal{R}(t) = \{(s, r) : s, r \in V,\ s \neq r\}$, but
\pkg{amore} also supports two-mode (bipartite) universes and
non-repeatable processes in which a dyad leaves the risk set once it has
fired (Section~\ref{sec:case-control}).

The relational event model treats the log as a realization of a marked
point process in which each dyad $(s,r)$ carries its own conditional
intensity
\begin{equation}
\lambda_{sr}(t \mid H_t)
  = \lambda_0(t)\,\exp\!\bigl\{\eta_{sr}(t \mid H_t)\bigr\},
  \qquad
\eta_{sr}(t \mid H_t) = \boldsymbol{\beta}^\top
   \mathbf{x}_{sr}(t \mid H_t),
\label{eq:intensity}
\end{equation}
where $H_t = \{e_j : t_j < t\}$ is the event history up to $t$,
$\mathbf{x}_{sr}(t \mid H_t)$ is a vector of \emph{statistics} that
summarize how the dyad's propensity to act depends on that history and on
exogenous attributes, $\boldsymbol{\beta}$ is the parameter of interest,
and $\lambda_0(t)$ is a baseline rate. We refer to
$\eta_{sr}(t \mid H_t)$ as the dyad's \emph{log-intensity} or
\emph{score}; it is the quantity each estimation backend models, linearly
or otherwise.

Conditioning on the observed event times, the baseline $\lambda_0$ drops
out and inference for $\boldsymbol{\beta}$ proceeds through the Cox-type
partial likelihood
\begin{equation}
\mathrm{PL}(\boldsymbol{\beta})
  = \prod_{i=1}^{M}
    \frac{\exp\!\bigl\{\eta_{s_i r_i}(t_i \mid H_{t_i})\bigr\}}
         {\sum_{(s,r) \in \mathcal{R}(t_i)}
            \exp\!\bigl\{\eta_{sr}(t_i \mid H_{t_i})\bigr\}},
\label{eq:partial-likelihood}
\end{equation}
in which the $i$-th factor is the probability that, among all dyads at
risk at $t_i$, the one that actually fired was $(s_i, r_i)$
\citep{butts2008, bianchi2024}. Equation~\eqref{eq:partial-likelihood} is
a product of softmax terms over the risk sets and is the estimand that
unifies the package: it is what the linear conditional logit and the
neural backend of Section~\ref{sec:framework} both maximize.

\subsection{Nested case-control sampling and the conditional-logistic
partial likelihood}
\label{sec:case-control}

The denominator of Equation~\eqref{eq:partial-likelihood} sums over the
entire risk set, which has $O(|V|^2)$ terms and must be recomputed (with
its history-dependent statistics) at every event. For networks of even
moderate size this is the computational bottleneck of REM estimation.
Nested case-control sampling removes it \citep{lerner2020}: at each event
time $t_i$ one retains the observed dyad (the \emph{case}) and draws a
small set of $k$ non-events (the \emph{controls}) uniformly at random from
$\mathcal{R}(t_i) \setminus \{(s_i, r_i)\}$. The case together with its
$k$ sampled controls forms a \emph{stratum} $\mathcal{S}_i$ of size
$k + 1$. Replacing the full risk set by $\mathcal{S}_i$ in
Equation~\eqref{eq:partial-likelihood} yields the sampled partial
likelihood
\begin{equation}
\mathrm{PL}_{\mathrm{cc}}(\boldsymbol{\beta})
  = \prod_{i=1}^{M}
    \frac{\exp\!\bigl\{\eta_{s_i r_i}(t_i \mid H_{t_i})\bigr\}}
         {\sum_{(s,r) \in \mathcal{S}_i}
            \exp\!\bigl\{\eta_{sr}(t_i \mid H_{t_i})\bigr\}},
\label{eq:cc-likelihood}
\end{equation}
which is exactly the partial likelihood of a \emph{stratified
conditional-logistic regression}: each stratum contributes one event
among $k+1$ candidates, and the contribution is the softmax of the
candidate scores. \citet{lerner2020} show that the resulting estimator is
consistent and that its efficiency relative to the full-risk-set fit is
governed by the number of controls $k$, with diminishing returns beyond a
handful. This is the design that \pkg{amore} adopts as its canonical
in-memory representation: a long table with one row per candidate,
carrying the event indicator, the stratum identifier, and the computed
statistics. We write $y_{ij} \in \{0,1\}$ for the event indicator of
candidate $j$ in stratum $\mathcal{S}_i$ (with $\sum_j y_{ij} = 1$) and
$\boldsymbol{\eta}_i = (\eta_{ij})_j$ for the stratum's score vector;
Equation~\eqref{eq:cc-likelihood} is then $\prod_i
\mathrm{softmax}(\boldsymbol{\eta}_i)_{j^\star(i)}$ evaluated at the event
position $j^\star(i)$.

\pkg{amore} builds this table with \texttt{sample\_non\_events()}, which
exposes the sampling \emph{scope} (drawing controls from all dyads, only
from already-appeared actors, or only along admissible citation edges),
the \emph{mode} (one- versus two-mode universes), and a \emph{risk}
policy that, for non-repeatable processes such as citation or species
invasion, removes a dyad from the risk set once it has fired and excludes
concurrent events from one another's control pools. Crucially, controls
read the \emph{true} network history $H_{t_i}$ when their statistics are
evaluated, but, being non-events, they never write to it; this read/write
split is what keeps the sampled statistics faithful to the full-data
intensity of Equation~\eqref{eq:intensity}.

\subsection{The degenerate case-1-control binomial reduction}
\label{sec:degenerate}

The conditional-logistic likelihood~\eqref{eq:cc-likelihood} admits
linear predictors only. To fit smooth, time-varying, and random effects
with the mature \pkg{mgcv} machinery, \pkg{amore} uses the specialization
of \citet{boschi2025gof, boschi2025rhem} to a single matched control,
$k = 1$. With exactly one case and one control per stratum, write
$\mathbf{x}^{\mathrm{ev}}_i$ and $\mathbf{x}^{\mathrm{nv}}_i$ for the
statistic vectors of the event and its matched non-event. The two-element
softmax in Equation~\eqref{eq:cc-likelihood} then collapses to a logistic
form in the \emph{difference} $\mathbf{x}^{\mathrm{ev}}_i -
\mathbf{x}^{\mathrm{nv}}_i$, which can be fitted as an ordinary binomial
regression with all responses equal to one and the linear predictor
entered through a paired event-minus-control contrast. \pkg{amore}'s wide
case-1-control representation materializes this directly, carrying for
each covariate the event value, the non-event value, and their difference
($\texttt{<cov>\_ev}$, $\texttt{<cov>\_nv}$, $\texttt{d\_<cov>}$), so that
a smooth term $f(\cdot)$ enters \pkg{mgcv} as the matched contrast
$f(\mathbf{x}^{\mathrm{ev}}_i) - f(\mathbf{x}^{\mathrm{nv}}_i)$.

We stress the precise relationship, because it is easy to overstate. The
degenerate binomial logistic regression of this subsection and the
conditional-logistic partial likelihood of
Equation~\eqref{eq:cc-likelihood} \emph{coincide} only in the single-control
case $k = 1$; they are algebraically related, not identical, and the
package does not constrain the conditional-logit and neural backends to
$k = 1$. Accordingly, throughout the paper we treat the smooth backend as
fitting the \emph{degenerate case-1-control binomial logistic regression}
of \citet{boschi2025gof}, distinct from---though kindred to---the risk-set
softmax partial likelihood that the linear and neural backends share.

\subsection{Endogenous statistics: families and variant axes}
\label{sec:statistics}

The statistics $\mathbf{x}_{sr}(t \mid H_t)$ are what make the model
\emph{relational}: they encode how prior events shape current propensities.
\pkg{amore} ships the endogenous-effect catalogue of
\citet{juozaitiene2024}, organized as a small number of structural
\emph{families} crossed with several \emph{variant axes}. The families are
the canonical sub-network configurations of directed interaction:
\emph{reciprocity} ($r$ has recently acted on $s$), \emph{transitivity}
and \emph{cyclic} closure (two-path and three-cycle configurations through
a common third actor), and \emph{sending} and \emph{receiving balance}
(shared-partner effects on the sender and receiver sides), together with
the elementary degree and recency statistics. Each family is then
instantiated along orthogonal variant axes that control \emph{how} past
events are aggregated:
\begin{itemize}
  \item \emph{counting} versus \emph{binary} --- the number of qualifying
    past events versus a $0/1$ indicator of their presence;
  \item \emph{exponential decay} and \emph{recency} (\texttt{time\_recent},
    \texttt{time\_first}) --- weighting past events by elapsed time, with a
    tunable half-life, rather than counting them equally;
  \item \emph{ordered} --- distinguishing the temporal order of the two
    legs of a closure configuration;
  \item \emph{interrupted} --- the central methodological variant of
    \citet{juozaitiene2024}, which switches a closure statistic off once an
    intervening event has broken the configuration, capturing the
    transience of social structure.
\end{itemize}
Crossing the families with these axes yields a catalogue of $74$ numeric
statistics, of which $66$ are computed through a C\texttt{++} fast path and
$8$ through a pure-R fallback, plus the list-valued
\texttt{sender\_receivers\_set} statistic used internally for
appearance-scoped sampling. The same engine,
\texttt{compute\_endogenous\_features()}, evaluates these statistics
post~hoc from timestamps alone, and (Section~\ref{sec:validation}) agrees
with the on-the-fly values recorded by the simulator to floating-point
precision. Exogenous covariates---actor attributes, dyadic distances, and
the like---enter $\eta_{sr}$ additively alongside the endogenous terms.

\subsection{Global covariates via a time-shifted partial likelihood}
\label{sec:global}

A purely time-varying covariate that is shared by every dyad at a given
instant---a calendar effect, weather, a system-wide shock---cancels from
the risk-set softmax of Equation~\eqref{eq:cc-likelihood}, because it
enters every candidate's score identically and so vanishes from the
ratio. Such \emph{global} covariates are therefore not identifiable from
the standard partial likelihood. \citet{lembo2025} resolve this with a
time-shifted construction: each candidate's global covariate is evaluated
not at the event time but at a dyad-specific shifted time determined by
its own history (for instance, the time since that dyad last acted), so
that the covariate no longer enters identically across the stratum and a
coefficient becomes estimable within the same conditional-logistic
framework. \pkg{amore} implements this scheme so that global covariates
can be carried through the identical case-control pipeline as endogenous
and exogenous statistics, with no change to the downstream fitter.

Together, Equations~\eqref{eq:intensity}--\eqref{eq:cc-likelihood}, the
$k=1$ reduction of Section~\ref{sec:degenerate}, the statistics catalogue
of Section~\ref{sec:statistics}, and the time-shift of
Section~\ref{sec:global} define the single case-control table that every
component of \pkg{amore} produces or consumes. The next section describes
how the package constructs that table.
